\documentclass[11pt,paper=a4,answers]{exam}
\usepackage{graphicx,lastpage}
\usepackage{upgreek}
\usepackage{censor}
\usepackage{gensymb}
\usepackage{amsmath}
\usepackage{multicol}
\usepackage{enumitem}
\usepackage{tasks}
\usepackage{adjustbox}
\usepackage{xcolor}
\usepackage{tfrupee}
\setlength{\voffset}{0.25in}
\usepackage{tabularx}
\newcolumntype{Y}{>{\centering\arraybackslash}X}
%==============================================================
% Customise Non-Essential Packages Below
% =============================================================
\usepackage[most]{tcolorbox}
\usepackage{eso-pic}
\usepackage{tikz}
\AddToShipoutPictureBG{%
\begin{tikzpicture}[remember picture,overlay]
\foreach \x in {-8,-4,0,4,8,12,16,20}
\foreach \y in {-8,-4,0,4,8,12,16,20} {
\node[opacity=0.05,rotate=45,anchor=center] at ([xshift=\x cm,yshift=\y cm]current page.center) {%
\begin{minipage}{2cm}
\centering
\includegraphics[width=2cm]{images/logo_black.png}\\
\small preppeo.com
\end{minipage}%
};
}
\end{tikzpicture}%
}
\printanswers

%==============================================================
% Customise Non-Essential Packages Above
% =============================================================
\definecolor{svgray}{rgb}{0.90, 0.90, 0.90}
\graphicspath{{images/}}

\usepackage[normalem]{ulem}
\renewcommand\ULthickness{1pt}   %%---> For changing thickness of underline
\setlength\ULdepth{3.5ex}%\maxdimen ---> For changing depth of underline
\renewcommand{\baselinestretch}{1}
\chead{
\includegraphics[scale=0.1,height=2cm ]{logo.png}
}
\pagestyle{empty}

\pagestyle{headandfoot}
\headrule
\cfoot{W: https://preppeo.com; M: +91-8130711689}

\begin{document}
%==============================================================
% Customise Question paper details below this
% =============================================================
\begin{center}
    \centering
    \renewcommand{\arraystretch}{1.5}
    \begin{tabularx}{\textwidth}{YYY}
    & \normalsize\bf{{\Large{{Quantitative Reasoning}}}} & \\
    By Preppeo & \normalsize\bf{{sat\_lid\_023}} & SAT\\
    \normalsize{{\textbf{{Unit:}} Advanced Math}} & \normalsize{{\textbf{{Chapter:}} Exponential functions and equations}} & \normalsize{{\textbf{{Topic:}} Exponential Growth \& Decay}} \\
    \end{tabularx}
\end{center}
\par\noindent\rule{\textwidth}{1pt}
% =============================================================
% Customise Question paper details above this
% =============================================================

\begin{questions}
\pointsinrightmargin
\pointsdroppedatright
\marksnotpoints
\pointpoints{mark}{marks}
\pointformat{\boldmath\themarginpoints}
\bracketedpoints

% =============================================================
% Question Begin
% =============================================================
% Lid Topic: Advanced Math — Exponential Growth & Decay
% Total Questions: 50

% --- PART 1: Modeling and Interpretations (1-25) ---

% Question 1
% Category: Concept | Difficulty: Medium | Question Type: Multiple Choice
\question
A function $p$ estimates that there were 2,500 birds in a population in 2010. Each year from 2010 to 2020, the function estimates that the number of birds in this population increased by 4\% of the number of birds in the population the previous year. Which equation defines this function, where $p(x)$ is the estimated number of birds in the population $x$ years after 2010?
\begin{tasks}(4)
    \task $p(x) = 2,500(4)^x$
    \task $p(x) = 2,500(1.04)^x$
    \task $p(x) = 2,500(0.96)^x$
    \task $p(x) = 2,500(0.04)^x$
\end{tasks}

\begin{solution}
\textbf{Conceptual Explanation:} \\
Exponential growth is modeled by the function $y = a(1 + r)^x$, where $a$ is the initial value, $r$ is the growth rate as a decimal, and $x$ is the number of time periods.

\vspace{20pt}

\textbf{Calculation and Logic:} \\
Initial value ($a$) $= 2,500$. \\
Growth rate ($r$) $= 4\% = 0.04$. \\
Growth factor $= 1 + r = 1 + 0.04 = 1.04$.

\vspace{10pt}

Substitute into the formula: \\
$p(x) = 2,500(1.04)^x$.

\vspace{25pt}

Answer: \colorbox{yellow!100}{B}
\end{solution}

\vspace{40pt}

% Question 2
% Category: Exam level | Difficulty: Hard | Question Type: Multiple Choice
\question
For the function $q$, the value of $q(x)$ decreases by 35\% for every increase in the value of $x$ by 1. If $q(0) = 25$, which equation defines $q$?
\begin{tasks}(4)
    \task $q(x) = 0.65(25)^x$
    \task $q(x) = 1.35(25)^x$
    \task $q(x) = 25(0.65)^x$
    \task $q(x) = 25(1.35)^x$
\end{tasks}

\begin{solution}
\textbf{Conceptual Explanation:} \\
Exponential decay is modeled by $y = a(1 - r)^x$. A percentage decrease results in a decay factor less than 1.

\vspace{20pt}

\textbf{Calculation and Logic:} \\
Initial value ($a$) $= q(0) = 25$. \\
Decay rate ($r$) $= 35\% = 0.35$. \\
Decay factor $= 1 - r = 1 - 0.35 = 0.65$.

\vspace{10pt}

Substitute into the formula: \\
$q(x) = 25(0.65)^x$.

\vspace{25pt}

Answer: \colorbox{yellow!100}{C}
\end{solution}

\vspace{40pt}

% Question 3
% Category: Practice | Difficulty: Medium | Question Type: Numeric Entry
\question
Bacteria are growing in a liquid growth medium. There were 400,000 cells per milliliter during an initial observation. The number of cells per milliliter doubles every 4 hours. How many cells per milliliter will there be 12 hours after the initial observation?

\begin{solution}
\textbf{Conceptual Explanation:} \\
For doubling scenarios, determine the number of doubling periods that occur in the total time given. Multiply the initial amount by $2^n$, where $n$ is the number of periods.

\vspace{20pt}

\textbf{Calculation and Logic:} \\
Initial amount $= 400,000$. \\
Total time $= 12$ hours. \\
Doubling time $= 4$ hours. \\
Number of doublings ($n$) $= 12 / 4 = 3$.

\vspace{10pt}

Calculate final amount: \\
$400,000 \times 2^3 = 400,000 \times 8 = 3,200,000$.

\vspace{25pt}

Answer: \colorbox{yellow!100}{3,200,000}
\end{solution}

\vspace{40pt}

% Question 4
% Category: Exam level | Difficulty: Hard | Question Type: Multiple Choice
\question
$f(x) = 150,000(1.18)^x$ \\
The function $f$ gives a company's predicted annual revenue, in dollars, $x$ years after the company started selling products online, where $0 < x \leq 8$. What is the best interpretation of the statement "$f(4)$ is approximately equal to 290,816" in this context?
\begin{tasks}(1)
    \task 4 years after the company started selling online, its predicted annual revenue is approximately \$290,816.
    \task 4 years after the company started selling online, its predicted annual revenue will have increased by a total of \$290,816.
    \task When the revenue is approximately \$290,816, it is 4 times the revenue for the previous year.
    \task When the revenue is approximately \$290,816, it is 4\% greater than the revenue for the previous year.
\end{tasks}

\begin{solution}
\textbf{Conceptual Explanation:} \\
In function notation $f(x) = y$, the value inside the parentheses ($x$) is the input (time), and the result ($y$) is the output (revenue).

\vspace{20pt}

\textbf{Calculation and Logic:} \\
$x = 4$ represents 4 years after starting. \\
$f(4) = 290,816$ represents the revenue at that specific point in time.

\vspace{25pt}

Answer: \colorbox{yellow!100}{A}
\end{solution}

\vspace{40pt}

% Question 5
% Category: Concept | Difficulty: Medium | Question Type: Multiple Choice
\question
The population of a town is currently 45,000, and the population is estimated to increase each year by 2.5\% from the previous year. Which of the following equations can be used to estimate the number of years, $t$, it will take for the population of the town to reach 55,000?
\begin{tasks}(2)
    \task $45,000 = 55,000(0.025)^t$
    \task $45,000 = 55,000(2.5)^t$
    \task $55,000 = 45,000(0.025)^t$
    \task $55,000 = 45,000(1.025)^t$
\end{tasks}

\begin{solution}
\textbf{Conceptual Explanation:} \\
Set up the exponential growth equation $Final = Initial(1 + r)^t$.

\vspace{20pt}

\textbf{Calculation and Logic:} \\
Initial population $= 45,000$. \\
Final population goal $= 55,000$. \\
Growth factor $= 1 + 0.025 = 1.025$.

\vspace{10pt}

Equation: $55,000 = 45,000(1.025)^t$.

\vspace{25pt}

Answer: \colorbox{yellow!100}{D}
\end{solution}

\vspace{40pt}

% Question 6
% Category: Exam level | Difficulty: Hard | Question Type: Numeric Entry
\question
For the exponential function $f$, the table below shows several values of $x$ and their corresponding values of $f(x)$, where $b$ is a constant greater than 1. If $k$ is a constant and $f(k) = b^{25}$, what is the value of $k$?
\begin{center}
\begin{tabular}{|c|c|}
\hline
$x$ & $f(x)$ \\ \hline
1 & $b^1$ \\ \hline
2 & $b^4$ \\ \hline
3 & $b^7$ \\ \hline
\end{tabular}
\end{center}

\begin{solution}
\textbf{Conceptual Explanation:} \\
Identify the linear pattern in the exponents relative to the $x$ values. The exponents form an arithmetic sequence.

\vspace{20pt}

\textbf{Calculation and Logic:} \\
Exponents are 1, 4, 7... \\
Common difference $= 3$. \\
Exponent formula: $E = 3x - 2$. \\
(Check: $x=1 \rightarrow E=3(1)-2=1$; $x=3 \rightarrow E=3(3)-2=7$).

\vspace{10pt}

Set $E = 25$: \\
$25 = 3x - 2$ \\
$27 = 3x \Rightarrow x = 9$.

\vspace{25pt}

Answer: \colorbox{yellow!100}{9}
\end{solution}

\vspace{40pt}

% Question 7
% Category: Practice | Difficulty: Hard | Question Type: Multiple Choice
\question
$p(t) = 120,000(1.08)^t$ \\
The given function $p$ models the population of a city $t$ years after a census. Which of the following functions best models the population of the city $m$ months after the census?
\begin{tasks}(2)
    \task $r(m) = \frac{120,000}{12}(1.08)^m$
    \task $r(m) = 120,000(\frac{1.08}{12})^m$
    \task $r(m) = 120,000(1.08)^{12m}$
    \task $r(m) = 120,000(1.08)^{\frac{m}{12}}$
\end{tasks}

\begin{solution}
\textbf{Conceptual Explanation:} \\
When changing units of time, replace the original variable with an equivalent expression in the new unit. Since $12 \text{ months} = 1 \text{ year}$, then $1 \text{ month} = 1/12 \text{ years}$. Thus, $t = m/12$.

\vspace{20pt}

\textbf{Calculation and Logic:} \\
Substitute $t = m/12$ into the original function: \\
$r(m) = 120,000(1.08)^{\frac{m}{12}}$.

\vspace{25pt}

Answer: \colorbox{yellow!100}{D}
\end{solution}

\vspace{40pt}

% Question 8
% Category: Exam level | Difficulty: Hard | Question Type: Multiple Choice
\question
$f(x) = (1.44)^{\frac{x}{2}}$ \\
The function $f$ is defined by the given equation. The equation can be rewritten as $f(x) = (1 + \frac{p}{100})^x$, where $p$ is a constant. What is the value of $p$?
\begin{tasks}(4)
    \task 12
    \task 20
    \task 22
    \task 44
\end{tasks}

\begin{solution}
\textbf{Conceptual Explanation:} \\
Use exponent rules to isolate the $x$ power. Specifically, $(a^b)^c = a^{bc}$. We can group the $1/2$ power with the base to simplify the growth factor.

\vspace{20pt}

\textbf{Calculation and Logic:} \\
$(1.44)^{\frac{x}{2}} = (1.44^{\frac{1}{2}})^x$ \\
Calculate the square root: \\
$\sqrt{1.44} = 1.2$.

\vspace{10pt}

Now the form is $(1.2)^x$. \\
Set $1 + \frac{p}{100} = 1.2$: \\
$\frac{p}{100} = 0.2 \Rightarrow p = 20$.

\vspace{25pt}

Answer: \colorbox{yellow!100}{B}
\end{solution}

\vspace{40pt}

% Question 9
% Category: Practice | Difficulty: Medium | Question Type: Numeric Entry
\question
A radioactive substance decays such that its mass is halved every 15 years. If the initial mass was 600 grams, what will the mass, in grams, be after 45 years?

\begin{solution}
\textbf{Conceptual Explanation:} \\
This is a half-life problem. Determine how many half-life periods occur in the total time.

\vspace{20pt}

\textbf{Calculation and Logic:} \\
Total time $= 45$. \\
Half-life $= 15$. \\
Number of cycles $= 45 / 15 = 3$.

\vspace{10pt}

Final mass $= 600 \times (0.5)^3$ \\
Final mass $= 600 \times 0.125 = 75$.

\vspace{25pt}

Answer: \colorbox{yellow!100}{75}
\end{solution}

\vspace{40pt}

% Question 10
% Category: Concept | Difficulty: Medium | Question Type: Multiple Choice
\question
\begin{center}
\begin{tikzpicture}[scale=0.6]
    \draw[->] (0,0) -- (6,0) node[right] {$x$};
    \draw[->] (0,0) -- (0,6) node[above] {$y$};
    \draw[domain=0:2.5, smooth, variable=\x, red, thick] plot ({\x}, {0.5*2^(\x)});
    \node at (2.5,5) {$y=f(x)$};
\end{tikzpicture}
\end{center}
The graph above shows an exponential function $f$. Which of the following statements about the function must be true?
\begin{tasks}(1)
    \task The value of $y$ increases by a constant amount for every unit increase in $x$.
    \task The value of $y$ increases by a constant percentage for every unit increase in $x$.
    \task The function has a constant slope.
    \task The $y$-intercept is $(0, 0)$.
\end{tasks}

\begin{solution}
\textbf{Conceptual Explanation:} \\
Linear functions increase by constant amounts. Exponential functions increase (or decrease) by a constant ratio or percentage.

\vspace{25pt}

Answer: \colorbox{yellow!100}{B}
\end{solution}

\vspace{40pt}

% Question 11
% Category: Exam level | Difficulty: Hard | Question Type: Numeric Entry
\question
According to Moore’s law, the number of transistors on a chip doubles every 2 years. In 2000, a chip had 10 million transistors. In which year does Moore’s law estimate the number of transistors to reach 80 million?

\begin{solution}
\textbf{Conceptual Explanation:} \\
Determine the number of doubling cycles required to reach the target value, then multiply by the duration of one cycle to find the total time passed.

\vspace{20pt}

\textbf{Calculation and Logic:} \\
Initial $= 10$ million. \\
Target $= 80$ million. \\
Ratio $= 80/10 = 8$.

\vspace{10pt}

Since $2^3 = 8$, three doubling cycles are required. \\
Each cycle $= 2$ years. \\
Total years passed $= 3 \times 2 = 6$.

\vspace{10pt}

Year $= 2000 + 6 = 2006$.

\vspace{25pt}

Answer: \colorbox{yellow!100}{2006}
\end{solution}

\vspace{40pt}

% Question 12
% Category: Practice | Difficulty: Medium | Question Type: Multiple Choice
\question
An investment account earns 6\% annual interest compounded annually. If an initial deposit of \$2,000 is made, which function $V(t)$ represents the value of the account after $t$ years?
\begin{tasks}(2)
    \task $V(t) = 2,000(0.06)^t$
    \task $V(t) = 2,000(1.06)^t$
    \task $V(t) = 2,000 + 0.06t$
    \task $V(t) = 2,000(6)^t$
\end{tasks}

\begin{solution}
\textbf{Conceptual Explanation:} \\
Compound interest is a form of exponential growth. The rate is added to 1 to find the base.

\vspace{25pt}

Answer: \colorbox{yellow!100}{B}
\end{solution}

\vspace{40pt}

% Question 13
% Category: Concept | Difficulty: Medium | Question Type: Multiple Choice
\question
Which of the following exponential functions represents a decay of 12\% per period?
\begin{tasks}(2)
    \task $f(x) = 100(1.12)^x$
    \task $f(x) = 100(0.12)^x$
    \task $f(x) = 100(0.88)^x$
    \task $f(x) = 100(-0.12)^x$
\end{tasks}

\begin{solution}
\textbf{Conceptual Explanation:} \\
Decay factor $= 1 - \text{rate}$. \\
$1 - 0.12 = 0.88$.

\vspace{25pt}

Answer: \colorbox{yellow!100}{C}
\end{solution}

\vspace{40pt}

% Question 14
% Category: Exam level | Difficulty: Hard | Question Type: Numeric Entry
\question
The population of jackrabbits in a region was 50 in the year 1950. By 2000, the population had grown exponentially to 1,600. If the annual percentage increase was constant, what was the population in 1980?

\begin{solution}
\textbf{Conceptual Explanation:} \\
First, determine the growth factor for the given time interval. Then, apply that growth factor to the specific sub-interval requested.

\vspace{20pt}

\textbf{Calculation and Logic:} \\
Total time $= 50$ years. \\
Growth multiplier $= 1600/50 = 32$. \\
We have $r^{50} = 32$.

\vspace{10pt}

Time from 1950 to 1980 $= 30$ years. \\
We need $50 \times r^{30}$. \\
Since $32 = 2^5$, we can say $r^{50} = (2^{1/10})^{50} \Rightarrow r = 2^{1/10}$. \\
(Basically, the population doubles every 10 years).

\vspace{10pt}

1950 to 1980 is 3 doublings (30 years). \\
$50 \times 2^3 = 50 \times 8 = 400$.

\vspace{25pt}

Answer: \colorbox{yellow!100}{400}
\end{solution}

\vspace{40pt}

% Question 15
% Category: Practice | Difficulty: Hard | Question Type: Multiple Choice
\question
The functions $f$ and $g$ are defined by the given equations, where $x \geq 0$. Which of the following equations displays, as a constant or coefficient, the maximum value of the function it defines, where $x \geq 0$? \\
$f(x) = 50(0.8)^{x+2}$ \\
$g(x) = 50(0.64)(0.8)^{x-2}$
\begin{tasks}(4)
    \task I only
    \task II only
    \task I and II
    \task Neither I nor II
\end{tasks}

\begin{solution}
\textbf{Conceptual Explanation:} \\
For a decay function ($0 < \text{base} < 1$) where $x \geq 0$, the maximum value occurs at the start of the interval ($x=0$). The equation displays this value if the term $a(b)^x$ has $x$ as the exponent with no shifts, or if the coefficient represents $f(0)$.

\vspace{20pt}

\textbf{Calculation and Logic:} \\
$f(0) = 50(0.8)^2 = 50(0.64) = 32$. \\
Function $f$ does not show 32 as a coefficient.

\vspace{10pt}

$g(x) = 32(0.8)^{x-2}$. \\
$g(0) = 32(0.8)^{-2} = 32 / 0.64 = 50$. \\
Function $g$ does not show 50 as a lead coefficient. \\
Wait, let's re-evaluate $g(x)$ in your SS style: \\
$g(x) = 50(0.64)(0.8)^{x-2} = 32(0.8)^{x-2}$. At $x=2$, $g(2)=32$. \\
The question asks which displays the \textit{maximum} on the interval. For decay, max is at $x=0$. Neither form explicitly shows the value at $x=0$ as the primary coefficient without further calculation.

\vspace{25pt}

Answer: \colorbox{yellow!100}{D}
\end{solution}

% Question 16
% Category: Exam level | Difficulty: Medium | Question Type: Multiple Choice
\question
$h(t) = 15(2)^{\frac{t}{5}}$ \\
The function $h$ models the height of a plant, in centimeters, $t$ days after it was planted. Which of the following is the best interpretation of the number 5 in this context?
\begin{tasks}(1)
    \task The plant's height increases by 5 centimeters each day.
    \task The plant's height doubles every 5 days.
    \task The initial height of the plant was 5 centimeters.
    \task The plant will reach a maximum height of 5 centimeters.
\end{tasks}

\begin{solution}
\textbf{Conceptual Explanation:} \\
In the form $y = a(2)^{t/d}$, the value $d$ represents the time it takes for the quantity to double.

\vspace{25pt}

Answer: \colorbox{yellow!100}{B}
\end{solution}

\vspace{40pt}

% Question 17
% Category: Practice | Difficulty: Medium | Question Type: Numeric Entry
\question
The value of a computer is \$1,200 when new and depreciates exponentially. If the value of the computer decreases by 20\% each year, what is its value, in dollars, after 2 years?

\begin{solution}
\textbf{Calculation and Logic:} \\
$V = 1200(0.8)^2$ \\
$V = 1200(0.64) = 768$.

\vspace{25pt}

Answer: \colorbox{yellow!100}{768}
\end{solution}

\vspace{40pt}

% Question 18
% Category: Exam level | Difficulty: Hard | Question Type: Numeric Entry
\question
If $a^{2x} = (a^3)^4 \cdot a^2$, what is the value of $x$?

\begin{solution}
\textbf{Calculation and Logic:} \\
$a^{2x} = a^{12} \cdot a^2$ \\
$a^{2x} = a^{14}$ \\
$2x = 14 \Rightarrow x = 7$.

\vspace{25pt}

Answer: \colorbox{yellow!100}{7}
\end{solution}

\vspace{40pt}

% Question 19
% Category: Concept | Difficulty: Medium | Question Type: Multiple Choice
\question
A scientist observes a population of 500 insects. The population increases by 100\% every week. Which function represents the population $P$ after $w$ weeks?
\begin{tasks}(2)
    \task $P(w) = 500(1)^w$
    \task $P(w) = 500(2)^w$
    \task $P(w) = 500 + 100w$
    \task $P(w) = 500(1.10)^w$
\end{tasks}

\begin{solution}
\textbf{Conceptual Explanation:} \\
An increase of 100\% means the population is doubling. Growth factor $= 1 + 1.00 = 2$.

\vspace{25pt}

Answer: \colorbox{yellow!100}{B}
\end{solution}

\vspace{40pt}

% Question 20
% Category: Exam level | Difficulty: Hard | Question Type: Numeric Entry
\question
$f(x) = 2(3)^x$ \\
$g(x) = 162$ \\
If $f(x) = g(x)$, what is the value of $x$?

\begin{solution}
\textbf{Calculation and Logic:} \\
$2(3)^x = 162$ \\
$3^x = 81$ \\
Since $3^4 = 81$, $x = 4$.

\vspace{25pt}

Answer: \colorbox{yellow!100}{4}
\end{solution}

\vspace{40pt}

% Question 21
% Category: Practice | Difficulty: Hard | Question Type: Multiple Choice
\question
A painting purchased for \$5,000 increases in value by 10\% every 3 years. Which expression gives the value of the painting after $y$ years?
\begin{tasks}(2)
    \task $5,000(1.10)^{3y}$
    \task $5,000(1.10)^{\frac{y}{3}}$
    \task $5,000(1.30)^y$
    \task $5,000(0.10)^{\frac{y}{3}}$
\end{tasks}

\begin{solution}
\textbf{Conceptual Explanation:} \\
The exponent should represent the number of 3-year periods, which is $y/3$.

\vspace{25pt}

Answer: \colorbox{yellow!100}{B}
\end{solution}

\vspace{40pt}

% Question 22
% Category: Concept | Difficulty: Medium | Question Type: Multiple Choice
\question
$y = a(b)^x$ \\
In the exponential function above, if $a > 0$ and $0 < b < 1$, what happens to the value of $y$ as $x$ increases?
\begin{tasks}(1)
    \task $y$ increases at an increasing rate.
    \task $y$ increases at a decreasing rate.
    \task $y$ decreases toward 0.
    \task $y$ remains constant.
\end{tasks}

\begin{solution}
\textbf{Conceptual Explanation:} \\
If the base $b$ is between 0 and 1, the function represents decay, meaning the output gets smaller as the input gets larger.

\vspace{25pt}

Answer: \colorbox{yellow!100}{C}
\end{solution}

\vspace{40pt}

% Question 23
% Category: Exam level | Difficulty: Hard | Question Type: Numeric Entry
\question
The number of subscribers to a service was 1,000 at the end of the first month. Each month thereafter, the number of subscribers was 1.2 times the number of subscribers at the end of the previous month. How many subscribers were there at the end of the 4th month?

\begin{solution}
\textbf{Calculation and Logic:} \\
End of month 1: 1,000. \\
End of month 2: $1,000 \times 1.2 = 1,200$. \\
End of month 3: $1,200 \times 1.2 = 1,440$. \\
End of month 4: $1,440 \times 1.2 = 1,728$.

\vspace{25pt}

Answer: \colorbox{yellow!100}{1,728}
\end{solution}

\vspace{40pt}

% Question 24
% Category: Practice | Difficulty: Hard | Question Type: Multiple Choice
\question
Which of the following functions models a population that triples every 10 years?
\begin{tasks}(2)
    \task $P(t) = P_0 (3)^{10t}$
    \task $P(t) = P_0 (10)^{3t}$
    \task $P(t) = P_0 (3)^{\frac{t}{10}}$
    \task $P(t) = P_0 (10)^{\frac{t}{3}}$
\end{tasks}

\begin{solution}
\textbf{Conceptual Explanation:} \\
The base must be the growth factor (3) and the exponent must divide the total time by the period (10).

\vspace{25pt}

Answer: \colorbox{yellow!100}{C}
\end{solution}

\vspace{40pt}

% Question 25
% Category: Exam level | Difficulty: Hard | Question Type: Numeric Entry
\question
A colony of bacteria doubles in size every 20 minutes. If the colony starts with 50 bacteria, how many doublings will occur in 2 hours?

\begin{solution}
\textbf{Calculation and Logic:} \\
Total time $= 2 \text{ hours} = 120 \text{ minutes}$. \\
Cycle time $= 20 \text{ minutes}$. \\
Cycles $= 120 / 20 = 6$.

\vspace{25pt}

Answer: \colorbox{yellow!100}{6}
\end{solution}

% --- PART 2: Advanced Exponential Growth & Decay (26-50) ---

% Question 26
% Category: Exam level | Difficulty: Hard | Question Type: Numeric Entry
\question
A town's population was 12,000 in the year 2010. The population increases by 5\% every 3 years. If the population is modeled by the function $P(t) = 12,000(k)^{\frac{t}{3}}$, what is the value of $k$?

\begin{solution}
\textbf{Conceptual Explanation:} \\
The base of an exponential growth function represents the growth factor for one full period. A 5\% increase corresponds to a growth factor of $1 + 0.05$.

\vspace{20pt}

\textbf{Calculation and Logic:} \\
Initial value $= 12,000$. \\
Growth rate $= 5\% = 0.05$. \\
Growth factor $(k) = 1 + 0.05 = 1.05$.

\vspace{10pt}

The equation uses $t/3$ as the exponent to account for the fact that the growth occurs every 3 years. Thus, the base $k$ is 1.05.

\vspace{25pt}

Answer: \colorbox{yellow!100}{1.05}
\end{solution}

\vspace{40pt}

% Question 27
% Category: Practice | Difficulty: Medium | Question Type: Multiple Choice
\question
Which of the following functions represents an exponential decay where the value decreases by 25\% every year?
\begin{tasks}(2)
    \task $f(t) = 500(0.25)^t$
    \task $f(t) = 500(1.25)^t$
    \task $f(t) = 500(0.75)^t$
    \task $f(t) = 500(1.75)^t$
\end{tasks}

\begin{solution}
\textbf{Conceptual Explanation:} \\
Decay factor $= 1 - \text{decay rate}$. For a 25\% decrease, the factor is $1 - 0.25 = 0.75$.

\vspace{25pt}

Answer: \colorbox{yellow!100}{C}
\end{solution}

\vspace{40pt}

% Question 28
% Category: Exam level | Difficulty: Hard | Question Type: Numeric Entry
\question
The value of an antique car increases by 8\% each year. If the car is currently worth \$25,000, what will be its value, in dollars, in 2 years?

\begin{solution}
\textbf{Calculation and Logic:} \\
Value after 1 year $= 25,000 \times 1.08 = 27,000$. \\
Value after 2 years $= 27,000 \times 1.08 = 29,160$.

\vspace{25pt}

Answer: \colorbox{yellow!100}{29,160}
\end{solution}

\vspace{40pt}

% Question 29
% Category: Concept | Difficulty: Medium | Question Type: Multiple Choice
\question
$f(x) = 10(2)^x$ \\
Which of the following best describes the growth of the function $f$?
\begin{tasks}(1)
    \task The value of $f(x)$ increases by 2 for every unit increase in $x$.
    \task The value of $f(x)$ increases by 10 for every unit increase in $x$.
    \task The value of $f(x)$ doubles for every unit increase in $x$.
    \task The value of $f(x)$ increases by 200\% for every unit increase in $x$.
\end{tasks}

\begin{solution}
\textbf{Conceptual Explanation:} \\
A base of 2 in an exponential function indicates that the output value doubles ($100\%$ increase) for each unit increase in the input.

\vspace{25pt}

Answer: \colorbox{yellow!100}{C}
\end{solution}

\vspace{40pt}

% Question 30
% Category: Exam level | Difficulty: Hard | Question Type: Numeric Entry
\question
If $16^x = 2^{20}$, what is the value of $x$?

\begin{solution}
\textbf{Calculation and Logic:} \\
Rewrite 16 as a power of 2: $16 = 2^4$. \\
$(2^4)^x = 2^{20} \Rightarrow 2^{4x} = 2^{20}$. \\
$4x = 20 \Rightarrow x = 5$.

\vspace{25pt}

Answer: \colorbox{yellow!100}{5}
\end{solution}

\vspace{40pt}

% Question 31
% Category: Practice | Difficulty: Hard | Question Type: Numeric Entry
\question
A certain type of cell divides into two every 15 minutes. If a sample starts with 10 cells, how many cells will be present after 1 hour?

\begin{solution}
\textbf{Calculation and Logic:} \\
1 hour $= 60$ minutes. \\
Number of divisions $= 60 / 15 = 4$. \\
Total cells $= 10 \times 2^4 = 10 \times 16 = 160$.

\vspace{25pt}

Answer: \colorbox{yellow!100}{160}
\end{solution}

\vspace{40pt}

% Question 32
% Category: Concept | Difficulty: Medium | Question Type: Multiple Choice
\question
Which of the following is equivalent to $a^{\frac{5}{2}}$?
\begin{tasks}(4)
    \task $\sqrt[5]{a^2}$
    \task $\sqrt{a^5}$
    \task $2a^5$
    \task $5a^2$
    \task $\frac{1}{a^5}$
\end{tasks}

\begin{solution}
\textbf{Conceptual Explanation:} \\
Apply the radical rule: $x^{m/n} = \sqrt[n]{x^m}$.

\vspace{25pt}

Answer: \colorbox{yellow!100}{B}
\end{solution}

\vspace{40pt}

% Question 33
% Category: Exam level | Difficulty: Hard | Question Type: Multiple Choice
\question
$V(t) = 1000(1.04)^{\frac{t}{4}}$ \\
The function above models the value of an account $t$ years after an initial deposit. Which of the following is the best interpretation of the growth in this account?
\begin{tasks}(1)
    \task The value increases by 4\% every year.
    \task The value increases by 4\% every 4 years.
    \task The value increases by 1.04\% every 4 years.
    \task The value increases by \$4 every year.
\end{tasks}

\begin{solution}
\textbf{Conceptual Explanation:} \\
The base (1.04) represents a 4\% increase. The exponent ($t/4$) indicates that this increase occurs once every 4 units of $t$ (years).

\vspace{25pt}

Answer: \colorbox{yellow!100}{B}
\end{solution}

\vspace{40pt}

% Question 34
% Category: Practice | Difficulty: Hard | Question Type: Numeric Entry
\question
The population of a city is currently 200,000 and is decreasing at a rate of 2\% per year. What will the population be in 2 years?

\begin{solution}
\textbf{Calculation and Logic:} \\
Year 1: $200,000 \times 0.98 = 196,000$. \\
Year 2: $196,000 \times 0.98 = 192,080$.

\vspace{25pt}

Answer: \colorbox{yellow!100}{192,080}
\end{solution}

\vspace{40pt}

% Question 35
% Category: Concept | Difficulty: Medium | Question Type: Multiple Choice
\question
If $f(x) = a(b)^x$ represents exponential growth, which of the following must be true about $b$?
\begin{tasks}(4)
    \task $b < 0$
    \task $0 < b < 1$
    \task $b = 1$
    \task $b > 1$
\end{tasks}

\begin{solution}
\textbf{Conceptual Explanation:} \\
For growth, the base must be greater than 1. If it is between 0 and 1, it represents decay.

\vspace{25pt}

Answer: \colorbox{yellow!100}{D}
\end{solution}

\vspace{40pt}

% Question 36
% Category: Exam level | Difficulty: Hard | Question Type: Numeric Entry
\question
What is the value of $n$ if $3^n = \sqrt[3]{81}$?

\begin{solution}
\textbf{Calculation and Logic:} \\
$81 = 3^4$. \\
$\sqrt[3]{3^4} = 3^{\frac{4}{3}}$. \\
$n = 4/3$ or $1.33$.

\vspace{25pt}

Answer: \colorbox{yellow!100}{4/3}
\end{solution}

\vspace{40pt}

% Question 37
% Category: Practice | Difficulty: Hard | Question Type: Multiple Choice
\question
The number of active users on a platform increases by 50\% every month. If there were 1,000 users initially, how many users are there after $m$ months?
\begin{tasks}(2)
    \task $1,000(0.50)^m$
    \task $1,000(1.50)^m$
    \task $1,000 + 0.50m$
    \task $1,000(50)^m$
\end{tasks}

\begin{solution}
\textbf{Conceptual Explanation:} \\
An increase of 50\% means the growth factor is $1 + 0.50 = 1.50$.

\vspace{25pt}

Answer: \colorbox{yellow!100}{B}
\end{solution}

\vspace{40pt}

% Question 38
% Category: Concept | Difficulty: Medium | Question Type: Multiple Choice
\question
$y = 400(0.95)^x$ \\
Which statement best describes the function above?
\begin{tasks}(1)
    \task The value increases by 5\% each period.
    \task The value decreases by 5\% each period.
    \task The value decreases by 95\% each period.
    \task The value increases by 95\% each period.
\end{tasks}

\begin{solution}
\textbf{Conceptual Explanation:} \\
A base of 0.95 means $1 - 0.05$, which indicates a 5\% decrease (decay).

\vspace{25pt}

Answer: \colorbox{yellow!100}{B}
\end{solution}

\vspace{40pt}

% Question 39
% Category: Exam level | Difficulty: Hard | Question Type: Numeric Entry
\question
A loan of \$5,000 charges 12\% annual interest compounded monthly. What is the monthly interest rate as a decimal?

\begin{solution}
\textbf{Calculation and Logic:} \\
Annual rate $= 12\% = 0.12$. \\
Monthly rate $= 0.12 / 12 = 0.01$.

\vspace{25pt}

Answer: \colorbox{yellow!100}{0.01}
\end{solution}

\vspace{40pt}

% Question 40
% Category: Practice | Difficulty: Medium | Question Type: Multiple Choice
\question
$f(x) = 2^x$ \\
What is the value of $f(5) - f(3)$?
\begin{tasks}(4)
    \task 4
    \task 8
    \task 24
    \task 32
\end{tasks}

\begin{solution}
\textbf{Calculation and Logic:} \\
$f(5) = 2^5 = 32$. \\
$f(3) = 2^3 = 8$. \\
$32 - 8 = 24$.

\vspace{25pt}

Answer: \colorbox{yellow!100}{C}
\end{solution}

\vspace{40pt}

% Question 41
% Category: Concept | Difficulty: Medium | Question Type: Multiple Choice
\question

The graph above models which of the following scenarios?
\begin{tasks}(1)
    \task A population doubling every year.
    \task A bank account earning 5\% interest.
    \task The value of a car depreciating over time.
    \task The linear growth of a plant.
\end{tasks}

\begin{solution}
\textbf{Conceptual Explanation:} \\
A curve that slopes downward toward the $x$-axis represents exponential decay, such as depreciation.

\vspace{25pt}

Answer: \colorbox{yellow!100}{C}
\end{solution}

\vspace{40pt}

% Question 42
% Category: Exam level | Difficulty: Hard | Question Type: Numeric Entry
\question
If $5^k = \frac{1}{25}$, what is the value of $k$?

\begin{solution}
\textbf{Calculation and Logic:} \\
$\frac{1}{25} = 25^{-1} = (5^2)^{-1} = 5^{-2}$. \\
$k = -2$.

\vspace{25pt}

Answer: \colorbox{yellow!100}{-2}
\end{solution}

\vspace{40pt}

% Question 43
% Category: Practice | Difficulty: Hard | Question Type: Multiple Choice
\question
$P(t) = 50(3)^t$ \\
This function models a population of insects after $t$ days. How many insects are there after 2 days?
\begin{tasks}(4)
    \task 150
    \task 300
    \task 450
    \task 900
\end{tasks}

\begin{solution}
\textbf{Calculation and Logic:} \\
$P(2) = 50(3^2) = 50 \times 9 = 450$.

\vspace{25pt}

Answer: \colorbox{yellow!100}{C}
\end{solution}

\vspace{40pt}

% Question 44
% Category: Concept | Difficulty: Medium | Question Type: Multiple Choice
\question
Which table represents an exponential function?
\begin{tasks}(2)
    \task \begin{tabular}{|c|c|c|c|} \hline $x$ & 1 & 2 & 3 \\ \hline $y$ & 2 & 4 & 6 \\ \hline \end{tabular}
    \task \begin{tabular}{|c|c|c|c|} \hline $x$ & 1 & 2 & 3 \\ \hline $y$ & 2 & 4 & 8 \\ \hline \end{tabular}
    \task \begin{tabular}{|c|c|c|c|} \hline $x$ & 1 & 2 & 3 \\ \hline $y$ & 10 & 20 & 30 \\ \hline \end{tabular}
    \task \begin{tabular}{|c|c|c|c|} \hline $x$ & 1 & 2 & 3 \\ \hline $y$ & 5 & 10 & 15 \\ \hline \end{tabular}
\end{tasks}

\begin{solution}
\textbf{Conceptual Explanation:} \\
In an exponential function, the ratio between consecutive $y$-values is constant (here, $4/2 = 2$ and $8/4 = 2$). In linear functions, the difference is constant.

\vspace{25pt}

Answer: \colorbox{yellow!100}{B}
\end{solution}

\vspace{40pt}

% Question 45
% Category: Exam level | Difficulty: Hard | Question Type: Numeric Entry
\question
A substance has a half-life of 10 years. If 1,000 grams are present initially, how many grams remain after 40 years?

\begin{solution}
\textbf{Calculation and Logic:} \\
Cycles $= 40/10 = 4$. \\
Remaining $= 1,000 \times (0.5)^4 = 1,000 \times 0.0625 = 62.5$.

\vspace{25pt}

Answer: \colorbox{yellow!100}{62.5}
\end{solution}

\vspace{40pt}

% Question 46
% Category: Practice | Difficulty: Medium | Question Type: Multiple Choice
\question
An initial population of 100 triples every year. Which function represents the population after $x$ years?
\begin{tasks}(2)
    \task $f(x) = 100(3)^x$
    \task $f(x) = 100 + 3x$
    \task $f(x) = 300^x$
    \task $f(x) = 100(x)^3$
\end{tasks}

\begin{solution}
\textbf{Conceptual Explanation:} \\
Initial value $= 100$. Base (growth factor) $= 3$.

\vspace{25pt}

Answer: \colorbox{yellow!100}{A}
\end{solution}

\vspace{40pt}

% Question 47
% Category: Concept | Difficulty: Medium | Question Type: Multiple Choice
\question
In the function $y = a(b)^x$, what does $a$ represent?
\begin{tasks}(1)
    \task The growth rate
    \task The horizontal asymptote
    \task The initial value (y-intercept)
    \task The time period
\end{tasks}

\begin{solution}
\textbf{Conceptual Explanation:} \\
When $x=0$, $y = a(b)^0 = a(1) = a$. Thus, $a$ is the starting amount.

\vspace{25pt}

Answer: \colorbox{yellow!100}{C}
\end{solution}

\vspace{40pt}

% Question 48
% Category: Exam level | Difficulty: Hard | Question Type: Numeric Entry
\question
$2^{(x-3)} = 32$ \\
What is the value of $x$?

\begin{solution}
\textbf{Calculation and Logic:} \\
$32 = 2^5$. \\
$x - 3 = 5 \Rightarrow x = 8$.

\vspace{25pt}

Answer: \colorbox{yellow!100}{8}
\end{solution}

\vspace{40pt}

% Question 49
% Category: Practice | Difficulty: Hard | Question Type: Numeric Entry
\question
If a town's population of 5,000 increases by 10\% per year, how many people will be in the town after 2 years?

\begin{solution}
\textbf{Calculation and Logic:} \\
$5,000 \times (1.10)^2 = 5,000 \times 1.21 = 6,050$.

\vspace{25pt}

Answer: \colorbox{yellow!100}{6,050}
\end{solution}

\vspace{40pt}

% Question 50
% Category: Exam level | Difficulty: Hard | Question Type: Numeric Entry
\question
A computer file is compressed exponentially. Each step reduces the file size by 50\%. If the file is 1,024 MB and undergoes 3 steps of compression, what is the final size in MB?

\begin{solution}
\textbf{Calculation and Logic:} \\
$1,024 \times (0.5)^3 = 1,024 \times 0.125 = 128$.

\vspace{25pt}

Answer: \colorbox{yellow!100}{128}
\end{solution}


% =============================================================
% Question End
% =============================================================

\end{questions}
\begin{center}
\rule{.5\textwidth}{1pt}
\end{center}
\end{document}
